\documentclass{amsart}
% For dealing with references we use the comment environment
\usepackage{verbatim}
\newenvironment{reference}{\comment}{\endcomment}
%\newenvironment{reference}{}{}
\newenvironment{slogan}{\comment}{\endcomment}
\newenvironment{history}{\comment}{\endcomment}

% For commutative diagrams we use Xy-pic
\usepackage[all]{xy}

% We use 2cell for 2-commutative diagrams.
\xyoption{2cell}
\UseAllTwocells

% We use multicol for the list of chapters between chapters
\usepackage{multicol}

% This is generally recommended for better output
\usepackage{lmodern}
\usepackage[T1]{fontenc}

% For cross-file-references

% Package for hypertext links:
\usepackage{hyperref}

% For any local file, say "hello.tex" you want to link to please

% Theorem environments.
%
\theoremstyle{plain}
\newtheorem{theorem}[subsection]{Theorem}
\newtheorem{proposition}[subsection]{Proposition}
\newtheorem{lemma}[subsection]{Lemma}

\theoremstyle{definition}
\newtheorem{definition}[subsection]{Definition}
\newtheorem{example}[subsection]{Example}
\newtheorem{exercise}[subsection]{Exercise}
\newtheorem{situation}[subsection]{Situation}

\theoremstyle{remark}
\newtheorem{remark}[subsection]{Remark}
\newtheorem{remarks}[subsection]{Remarks}

\numberwithin{equation}{subsection}

% Macros
%
\def\lim{\mathop{\mathrm{lim}}\nolimits}
\def\colim{\mathop{\mathrm{colim}}\nolimits}
\def\Spec{\mathop{\mathrm{Spec}}}
\def\Hom{\mathop{\mathrm{Hom}}\nolimits}
\def\Ext{\mathop{\mathrm{Ext}}\nolimits}
\def\SheafHom{\mathop{\mathcal{H}\!\mathit{om}}\nolimits}
\def\SheafExt{\mathop{\mathcal{E}\!\mathit{xt}}\nolimits}
\def\Sch{\mathit{Sch}}
\def\Mor{\mathop{\mathrm{Mor}}\nolimits}
\def\Ob{\mathop{\mathrm{Ob}}\nolimits}
\def\Sh{\mathop{\mathit{Sh}}\nolimits}
\def\NL{\mathop{N\!L}\nolimits}
\def\CH{\mathop{\mathrm{CH}}\nolimits}
\def\proetale{{pro\text{-}\acute{e}tale}}
\def\etale{{\acute{e}tale}}
\def\QCoh{\mathit{QCoh}}
\def\Ker{\mathop{\mathrm{Ker}}}
\def\Im{\mathop{\mathrm{Im}}}
\def\Coker{\mathop{\mathrm{Coker}}}
\def\Coim{\mathop{\mathrm{Coim}}}

% Boxtimes
%
\DeclareMathSymbol{\boxtimes}{\mathbin}{AMSa}{"02}

%
% Macros for moduli stacks/spaces
%
\def\QCohstack{\mathcal{QC}\!\mathit{oh}}
\def\Cohstack{\mathcal{C}\!\mathit{oh}}
\def\Spacesstack{\mathcal{S}\!\mathit{paces}}
\def\Quotfunctor{\mathrm{Quot}}
\def\Hilbfunctor{\mathrm{Hilb}}
\def\Curvesstack{\mathcal{C}\!\mathit{urves}}
\def\Polarizedstack{\mathcal{P}\!\mathit{olarized}}
\def\Complexesstack{\mathcal{C}\!\mathit{omplexes}}
% \Pic is the operator that assigns to X its picard group, usage \Pic(X)
% \Picardstack_{X/B} denotes the Picard stack of X over B
% \Picardfunctor_{X/B} denotes the Picard functor of X over B
\def\Pic{\mathop{\mathrm{Pic}}\nolimits}
\def\Picardstack{\mathcal{P}\!\mathit{ic}}
\def\Picardfunctor{\mathrm{Pic}}
\def\Deformationcategory{\mathcal{D}\!\mathit{ef}}

\setlength{\emergencystretch}{3em}
\begin{document}
\title[Stacks Project, Tag 0DQ1]{416 --- Smooth local maps preserve numbers of geometric branches}
\maketitle
\noindent Copyright (C) 2005--2025 Johan de Jong. Stacks Project authors. Source: \href{https://stacks.math.columbia.edu/tag/0DQ1}{Tag 0DQ1}.

\noindent Editorial difficulty note (not author text). Difficulty H2: The author compares minimal primes after henselization and strict henselization using flat going down. Normalization and integral tensor products prove uniqueness of a prime above each base branch, with a separate residue-field argument required for the non-geometric assertion..

\noindent Editorial provenance note (not author text). History: excerpt from revision \texttt{a0444\allowbreak 6e57e\allowbreak c1fbc\allowbreak 252a8\allowbreak 71afc\allowbreak ec775\allowbreak 2fb28\allowbreak 07b14}, more-algebra.tex, lines 31595--31686. Reference presentation was adapted; original-author context and core support blocks were retained or added. The mathematical wording is unchanged. Licensed under GNU FDL 1.2 or later; no invariant sections or cover texts. See ../licenses/COPYING. Context blocks reproduce author definitions and supporting results; they are not additional counted problems.

\begin{lemma}
\label{lemma-invariance-number-branches-smooth}
Let $A \to B$ be a local homomorphism of local rings which
is the localization of a smooth ring map.
\begin{enumerate}
\item The number of geometric branches of $A$ is equal to the number
of geometric branches of $B$.
\item If $A \to B$ induces a purely inseparable extension of residue fields,
then the number of branches of $A$ is the number of branches of $B$.
\end{enumerate}
\end{lemma}

\begin{proof}
We will use that smooth ring maps are flat
(Algebra, Lemma \href{https://stacks.math.columbia.edu/tag/00TA}{lemma smooth syntomic (Tag 00TA)}),
that localizations are flat (Algebra, Lemma
\href{https://stacks.math.columbia.edu/tag/00HT}{lemma flat localization (Tag 00HT)}),
that compositions of flat ring maps are flat
(Algebra, Lemma \href{https://stacks.math.columbia.edu/tag/00HC}{lemma composition flat (Tag 00HC)}),
that base change of a flat ring map is flat
(Algebra, Lemma \href{https://stacks.math.columbia.edu/tag/00HI}{lemma flat base change (Tag 00HI)}),
that flat local homomorphisms are faithfully flat
(Algebra, Lemma \href{https://stacks.math.columbia.edu/tag/00HR}{lemma local flat ff (Tag 00HR)}),
that (strict) henselization is flat
(Lemma \href{https://stacks.math.columbia.edu/tag/07QM}{lemma dumb properties henselization (Tag 07QM)}), and
Going down for flat ring maps
(Algebra, Lemma \href{https://stacks.math.columbia.edu/tag/00HS}{lemma flat going down (Tag 00HS)}).

\medskip\noindent
Proof of (2). Let $A^h$, $B^h$ be the henselizations of $A$, $B$. Then $B^h$
is the henselization of $A^h \otimes_A B$ at the unique
maximal ideal lying over $\mathfrak m_B$, see
Algebra, Lemma \href{https://stacks.math.columbia.edu/tag/08HU}{lemma henselian functorial improve (Tag 08HU)}.
Thus we may and do assume $A$ is henselian.
Since $A \to B \to B^h$ is flat, every minimal prime of $B^h$
lies over a minimal prime of $A$ and since $A \to B^h$ is faithfully flat,
every minimal prime of $A$ does lie under a minimal prime of $B^h$;
in both cases use going down for flat ring maps.
Therefore it suffices to show that given a minimal prime
$\mathfrak p \subset A$, there is at most one minimal prime
of $B^h$ lying over $\mathfrak p$.
After replacing $A$ by $A/\mathfrak p$ and $B$ by $B/\mathfrak p B$
we may assume that $A$ is a domain; the $A$ is still henselian by
Algebra, Lemma \href{https://stacks.math.columbia.edu/tag/05WQ}{lemma quotient henselization (Tag 05WQ)}.
By Lemma \href{https://stacks.math.columbia.edu/tag/0BQ0}{lemma unibranch (Tag 0BQ0)} we see that the integral closure
$A'$ of $A$ in its field of fractions is a local domain.
Of course $A'$ is a normal domain. By
Algebra, Lemma \href{https://stacks.math.columbia.edu/tag/033C}{lemma normal goes up (Tag 033C)}
we see that $A' \otimes_A B^h$ is a normal ring
(the lemma just gives it for $A' \otimes_A B$, to go up to
$A' \otimes_A B^h$ use that $B^h$ is a colimit of \'etale
$B$-algebras and use Algebra, Lemma \href{https://stacks.math.columbia.edu/tag/037D}{lemma colimit normal ring (Tag 037D)}).
By Algebra, Lemma \href{https://stacks.math.columbia.edu/tag/092Y}{lemma local tensor with integral (Tag 092Y)}
we see that $A' \otimes_A B^h$ is local (this is where we
use the assumption on the residue fields of $A$ and $B$).
Hence $A' \otimes_A B^h$ is a local normal ring, hence a local domain.
Since $B^h \subset A' \otimes_A B^h$ by flatness of $A \to B^h$
we conclude that $B^h$ is a domain as desired.

\medskip\noindent
Proof of (1). Let $A^{sh}$, $B^{sh}$ be strict henselizations
of $A$, $B$. Then $B^{sh}$ is a strict henselization of
$A^h \otimes_A B$ at a maximal ideal lying over $\mathfrak m_B$
and $\mathfrak m_{A^h}$, see
Algebra, Lemma \href{https://stacks.math.columbia.edu/tag/08HV}{lemma strictly henselian functorial improve (Tag 08HV)}.
Thus we may and do assume $A$ is strictly henselian.
Since $A \to B \to B^{sh}$ is flat, every minimal prime of $B^{sh}$
lies over a minimal prime of $A$ and since $A \to B^{sh}$ is faithfully flat,
every minimal prime of $A$ does lie under a minimal prime of $B^{sh}$;
in both cases use going down for flat ring maps.
Therefore it suffices to show that given a minimal prime
$\mathfrak p \subset A$, there is at most one minimal prime
of $B^{sh}$ lying over $\mathfrak p$.
After replacing $A$ by $A/\mathfrak p$ and $B$ by $B/\mathfrak p B$
we may assume that $A$ is a domain; then $A$ is still strictly henselian by
Algebra, Lemma \href{https://stacks.math.columbia.edu/tag/05WS}{lemma quotient strict henselization (Tag 05WS)}.
By Lemma \href{https://stacks.math.columbia.edu/tag/06DM}{lemma geometrically unibranch (Tag 06DM)} we see that the integral closure
$A'$ of $A$ in its field of fractions is a local domain whose residue
field is a purely inseparable extension of the residue field of $A$.
Of course $A'$ is a normal domain. By
Algebra, Lemma \href{https://stacks.math.columbia.edu/tag/033C}{lemma normal goes up (Tag 033C)}
we see that $A' \otimes_A B^{sh}$ is a normal ring
(the lemma just gives it for $A' \otimes_A B$, to go up to
$A' \otimes_A B^{sh}$ use that $B^{sh}$ is a colimit of \'etale
$B$-algebras and use Algebra, Lemma \href{https://stacks.math.columbia.edu/tag/037D}{lemma colimit normal ring (Tag 037D)}).
By Algebra, Lemma \href{https://stacks.math.columbia.edu/tag/092Y}{lemma local tensor with integral (Tag 092Y)}
we see that $A' \otimes_A B^{sh}$ is local (since $A \subset A'$
induces a purely inseparable residue field extension).
Hence $A' \otimes_A B^{sh}$ is a local normal ring, hence a local domain.
Since $B^{sh} \subset A' \otimes_A B^{sh}$ by flatness of $A \to B^{sh}$
we conclude that $B^{sh}$ is a domain as desired.
\end{proof}
\end{document}
